% =========================================================
% Chapter: Formalizing the Number Systems
% =========================================================
\chapter{Formalizing the Number Systems}
\label{chap:formalizing-number-systems}

\begin{topicbox}{Why Formalization Is Needed}
\begin{exposition}
Ordinary arithmetic is one of the most familiar parts of mathematics. We
learn to count, add, subtract, multiply, and divide long before we ask what
sort of objects numbers are or why the usual rules are justified. For
foundational work, however, intuition is not enough. The basic number systems
must be specified by precise axioms and constructions, so that their
operations and order relations are not merely assumed but accounted for.
\end{exposition}
\end{topicbox}

\begin{topicbox}{The Natural-Number Structure}
\begin{exposition}
The first place where this becomes visible is the natural numbers. A
foundational treatment cannot simply point to the counting numbers and begin
calculating. It has to identify a first element, a successor operation, and
the induction principle that makes recursive definitions and proofs possible.
Peano's role was to isolate this structure: a first element, a successor map,
conditions saying that successor is injective and never returns to the first
element, and an induction principle strong enough to characterize the system.
\end{exposition}
\end{topicbox}

\begin{topicbox}{Structural Characterization}
\begin{exposition}
Dedekind's work emphasized the same structural point of view. What matters is
not the particular names attached to the elements, but the pattern determined
by a distinguished first element, successor, recursion, and induction. This
perspective underlies the treatment of Peano systems earlier in the volume,
and Dedekind's name also sits behind the cut-based construction of the real
numbers developed later in the usual foundational tradition.
\end{exposition}
\end{topicbox}

\begin{topicbox}{Later Number Systems}
\begin{exposition}
Once the natural numbers have been fixed, the later systems require their own
explicit constructions. The integers are built so that subtraction becomes an
internal operation. The rational numbers are built so that division by
nonzero quantities is represented inside the system. The real numbers are
then introduced to address the completeness failures of the rationals. In
each case, the construction supplies the operations, the order, and the
structural properties needed for analysis, rather than treating them as
background facts.
\end{exposition}
\end{topicbox}
